% Provenance (never printed): see sources/outline_sources.md (Year 3:
% partition coefficients; batch and fractional melting; crystallisation
% models; reading a spider diagram quantitatively). Brief:
% sources/book5/BRIEFS.md, ch. 12. Curves from figdata/bachelor-3/
% melting_models.py. Exercise and problem values in the Book 5 scratch
% calculator calc/ch12.py. Aldren Island and its lavas are invented.
\chapter{ट्रेस-तत्त्व मॉडलिंग}\label{ch:b3:trace-element-modelling}

एक ही द्वीप पर कुछ किलोमीटर की दूरी पर दो बेसाल्ट उद्गारित हुए। इनमें से एक में
दूसरे की तुलना में डेढ़ गुने से अधिक लैंथेनम और लगभग चार गुना कम इटरबियम होता है।
क्या ये दो अलग-अलग प्रावारों से गले, या एक ही प्रावार से
अलग-अलग मात्रा में गलकर, या अलग-अलग गहराई पर? मुख्य तत्त्व
मुश्किल से ही कुछ बताते हैं: दोनों ही बेसाल्ट हैं। अनुरेखी तत्त्व, जो खनिजों का अनुसरण करते हैं,
उन्हें बदले बिना, यह अभिलेख रखते हैं कि कितना गलित बना, वह अपने स्रोत से कैसे निकला
और कौन से खनिज पीछे रह गए। यह अध्याय उस अभिलेख को
समीकरणों में लिखता है और उसे फिर से पढ़ता है।

\begin{recall}
खनिज--गलित विभाजन गुणांक $D$, समग्र विभाजन गुणांक,
संगत और असंगत तत्त्व, दुर्लभ मृदा तत्त्व, CI
कॉन्ड्राइट और आद्य प्रावार, प्रसामान्यीकृत आरेख और यूरोपियम
विसंगति, MORB और OIB पुस्तक 4 (स्नातक वर्ष 2) में हैं, उसके अध्याय
पृथ्वी के तत्त्वों पर है। पेरिडोटाइट, लर्ज़ोलाइट, आंशिक गलन और गलन-अंश $F$
उसके उस अध्याय में हैं जो प्रावार के गलन पर है; साम्य क्रिस्टलन और
भिन्नात्मक क्रिस्टलन, रैले प्रभाजन नियम, मिश्रण और
आत्मसातीकरण--भिन्नात्मक क्रिस्टलन (AFC), जो मैग्मीय विभेदन पर उसके अध्याय में है।
वहाँ यह सिद्ध हुआ था कि क्रिस्टलन के दौरान दो अत्यंत असंगत
तत्त्वों का अनुपात स्थिर रहता है।
\end{recall}

\section{गलती चट्टान का लेखा-जोखा}\label{sec:b3:trace-element-modelling:balance}

द्रव्यमान 1 की एक चट्टान गलती है। गलन-अंश $F$ पर वह द्रव्यमान $F$ का
गलित बनाती है और द्रव्यमान $1 - F$ का ठोस अवशेष रखती है। एक अनुरेखी तत्त्व, जो स्रोत में
सांद्रता $C_0$ पर उपस्थित है, दोनों के बीच बँट जाता है; गलित में उसकी सांद्रता
गलित में उसकी सांद्रता $C_L$ है, अवशेष में $C_S$। इस अध्याय का हर मॉडल एक ही
संतुलन का एक रूप है, $C_0 = FC_L + (1 - F)C_S$, समूचे गलित और अवशेष के लिए,
और $C_S$ को $C_L$ से जोड़ने वाला एक नियम: किसी क्रिस्टल और
गलित के संपर्क पर, $C_{\text{खनिज}} = D\,C_L$। मॉडल इस बात में भिन्न होते हैं कि क्या किससे
छूता है और कब।

प्रावार के खनिजों के विभाजन गुणांक बहुत भिन्न होते हैं। एक संकलन में
लैंथेनम का $D = 0.058$ क्लिनोपाइरोक्सीन में है और $0.001$
गार्नेट में और ओलिवीन, ऑर्थोपाइरोक्सीन और स्पिनेल में; इटरबियम का $D = 0.58$ क्लिनोपाइरोक्सीन में
है, गार्नेट में $4.5$ और अन्य तीनों में $0.16$। गार्नेट भारी दुर्लभ मृदाओं को
ग्रहण करता है और हल्कों को अस्वीकार करता है।
% ledger: budahn91:d-ree

\begin{definition}[अ-मोडल गलन]\label{def:b3:trace-element-modelling:nonmodal}
\emph{अ-मोडल गलन}\index{अ-मोडल गलन} में स्रोत के खनिज गलित में अनुपात $p_i$ में प्रवेश करते हैं, जो
स्रोत में उनके द्रव्यमान भिन्नों $x_i$ से भिन्न होता है। गलित तब, आरंभ में,
स्रोत का समग्र विभाजन गुणांक $D_0 = \sum_i x_iD_i$ देखता है, और
जैसे-जैसे वह बढ़ता है, गलन-भारित गुणांक
$P = \sum_i p_iD_i$ देखता है।
\end{definition}

गलन शायद ही कभी मोडल होता है। स्पिनेल लर्ज़ोलाइट के एक हल किए उदाहरण में, जिसमें
\qty{58}{\percent} ओलिवीन, \qty{27}{\percent} ऑर्थोपाइरोक्सीन,
\qty{12}{\percent} क्लिनोपाइरोक्सीन और \qty{3}{\percent} स्पिनेल हैं, खनिज
20, 25, 45 और \qty{10}{\percent} के अनुपात में गलित में प्रवेश करते हैं:
क्लिनोपाइरोक्सीन अपने हिस्से से चार गुना तेज़ गलता है, ओलिवीन
अपने हिस्से से तीन गुना धीमे।
% ledger: white13:ex73

\begin{proposition}[अवशेष का समग्र विभाजन गुणांक]
\label{prop:b3:trace-element-modelling:residue}
स्थिर $p_i$ पर \omterm{def:b3:trace-element-modelling:nonmodal}{अ-मोडल गलन} में, गलन-अंश $F$ पर अवशेष का
समग्र विभाजन गुणांक इस प्रकार होता है:
\[
\bar D(F) = \frac{D_0 - PF}{1 - F},
\]
यह सूत्र तब तक मान्य है जब तक हर खनिज शेष है, अर्थात $F < x_i/p_i$, प्रत्येक $i$ के लिए।
\end{proposition}

\begin{proof}
स्रोत के प्रति इकाई द्रव्यमान पर, खनिज $i$ का, $F$ पर बचा द्रव्यमान $x_i - p_iF$ है,
और अवशेष का द्रव्यमान $1 - F$ है। उसका समग्र गुणांक द्रव्यमान-भारित
माध्य है, $\sum_i (x_i - p_iF)D_i/(1 - F) = (D_0 - PF)/(1 - F)$। कोई खनिज तब
समाप्त होता है जब $x_i - p_iF = 0$।
\end{proof}

ऊपर के खनिज गुणांकों और हल किए उदाहरण के स्पिनेल लर्ज़ोलाइट के साथ, तथा
एक गार्नेट लर्ज़ोलाइट, जो \qty{5}{GPa} और
\qty{1400}{\celsius} पर प्रयोगशाला में रखा गया, जिसमें \qty{54.1}{\percent} ओलिवीन,
\qty{9.8}{\percent} ऑर्थोपाइरोक्सीन, \qty{19.5}{\percent} क्लिनोपाइरोक्सीन और
\qty{16.7}{\percent} गार्नेट (द्रव्यमान के अनुसार) है, दोनों के समग्र गुणांक नीचे दिए गए हैं।
% ledger: white13:ex73, rfr25:gkr-5gpa, budahn91:d-ree
\begin{center}
\small
\begin{tabular}{lccc}
\hline
स्रोत & $D_{\mathrm{La}}$ & $D_{\mathrm{Sm}}$ & $D_{\mathrm{Yb}}$ \\
\hline
स्पिनेल लर्ज़ोलाइट & 0.0078 & 0.062 & 0.21 \\
गार्नेट लर्ज़ोलाइट & 0.012 & 0.107 & 0.97 \\
\hline
\end{tabular}
\end{center}
दोनों में लैंथेनम अत्यंत असंगत है। इटरबियम स्पिनेल लर्ज़ोलाइट में मध्यम रूप से असंगत है
और गार्नेट लर्ज़ोलाइट में लगभग उदासीन है। दोनों
स्रोतों के बीच का अंतर लगभग पूरी तरह भारी दुर्लभ
मृदाओं में है।

\section{बैच गलन}\label{sec:b3:trace-element-modelling:batch}

\begin{theorem}[बैच गलन]\label{thm:b3:trace-element-modelling:batch}
\emph{बैच गलन}\index{बैच गलन} में गलित पूरे अवशेष के साथ साम्य में
बना रहता है जब तक वह एक ही बैच में अलग नहीं होता। स्थिर समग्र
विभाजन गुणांक $D$ के साथ, गलन-अंश $F$ पर गलित और अवशेष
के लिए ये संबंध होते हैं:
\[
\frac{C_L}{C_0} = \frac{1}{D + F(1 - D)}, \qquad
\frac{C_S}{C_0} = \frac{D}{D + F(1 - D)}.
\]
\end{theorem}

\begin{proof}
संतुलन में अवशेष का हर क्रिस्टल $D_iC_L$ धारण करता है, इसलिए अवशेष
औसतन $C_S = DC_L$ धारण करता है। संतुलन $C_0 = FC_L + (1 - F)DC_L$ से
$C_L$ प्राप्त होता है, और $C_S$ उसी से निकलता है।
\end{proof}

यह सूत्र पुस्तक 4 (स्नातक वर्ष 2) के साम्य क्रिस्टलन नियम को उलटी ओर
पढ़ने का परिणाम है: दोनों में $F$ एक बंद तंत्र में द्रव का अंश है। इसकी दो सीमाएँ जानना उपयोगी है।
जब $F \to 0$ होता है, तो $C_L \to C_0/D$ होता है: गलित की पहली बूँद
किसी तत्त्व को अधिक-से-अधिक $1/D$ गुणक से समृद्ध कर सकती है। जब $D \to 0$ होता है, तो $C_L \to
C_0/F$ होता है: पूर्णतः असंगत तत्त्व गलित में तनु हो जाता है, और उसका
समृद्धीकरण केवल $F$ मापता है। संगत तत्त्व के लिए, $D > 1$, गलित
स्रोत से निर्धन होता है, और $F$ जितना छोटा हो, उतना ही अधिक निर्धन होता है।

\begin{proposition}[अ-मोडल बैच गलन]\label{prop:b3:trace-element-modelling:nonmodal}
\omterm{thm:b3:trace-element-modelling:batch}{बैच गलन} में गलन-अनुपात $p_i$ के साथ, जैसा \cref{def:b3:trace-element-modelling:nonmodal} में है,
\[
\frac{C_L}{C_0} = \frac{1}{D_0 + F(1 - P)}.
\]
\end{proposition}

\begin{proof}
अवशेष $\bar D(F)\,C_L$ धारण करता है (\cref{prop:b3:trace-element-modelling:residue} देखिए),
इसलिए $C_0 = FC_L + (1 - F)\bar D(F)C_L = C_L(F + D_0 - PF)$।
\end{proof}

जब $P$ $D_0$ से अधिक हो, जैसा हल किए गए उदाहरण में लैंथेनम के लिए होता है, जहाँ $D_0 =
0.0074$ और $P = 0.025$ हैं, तो अवशेष अपना सर्वोत्तम आश्रय खनिज, क्लिनोपाइरॉक्सीन, बाकी से तेज़ी से खोने लगता है।
तब गलित मोडल गलन के पूर्वानुमान से कुछ अधिक समृद्ध होता है,
$F = 0.1$ पर लैंथेनम के लिए \qty{2}{\percent} से कम (अभ्यास~6 देखिए), परंतु
अवशेष में क्लिनोपाइरॉक्सीन $F = 0.12/0.45 \approx 0.27$ पर समाप्त हो जाता है, और उससे आगे
खनिज और गुणांक बदल जाते हैं।
% ledger: white13:ex73

\section{भिन्नात्मक और गतिक गलन}\label{sec:b3:trace-element-modelling:fractional}

पेरिडोटाइट में गलित कणों के किनारों के साथ एकत्र होता है, और जैसे ही वह
जुड़ जाता है, निकल सकता है। \omterm{thm:b3:trace-element-modelling:batch}{बैच गलन} का विपरीत अंत्य-सदस्य हर बूँद को उसके
बनते ही हटा देता है।

\begin{theorem}[भिन्नात्मक गलन]\label{thm:b3:trace-element-modelling:fractional}
\emph{भिन्नात्मक गलन}\index{भिन्नात्मक गलन} में गलित का हर वृद्धिशील भाग
केवल अपने बनने के क्षण पर अवशेष के साथ साम्य में होता है, और फिर
उसे छोड़ देता है। स्थिर $D$ के साथ, गलित का कुल द्रव्यमान $F$ निकाला जा चुका हो, तब
अवशेष और उस क्षण बनने वाले गलित के लिए ये संबंध होते हैं:
\[
\frac{C_S}{C_0} = (1 - F)^{1/D - 1}, \qquad
\frac{C_L}{C_0} = \frac{1}{D}(1 - F)^{1/D - 1}.
\]
\end{theorem}

\begin{proof}
अवशेष का द्रव्यमान $1 - F$ और सांद्रता $C_S$ है। जब एक द्रव्यमान
$\mathrm dF$ गलित $C_L = C_S/D$ पर निकलता है, तो अवशेष में तत्त्व
$\mathrm d[(1 - F)C_S] = -(C_S/D)\,\mathrm dF$ से घट जाता है। विस्तार करने पर,
$(1 - F)\,\mathrm dC_S = C_S(1 - 1/D)\,\mathrm dF$ लिखा जा सकता है, इसलिए
$\mathrm dC_S/C_S = (1/D - 1)\,\mathrm d[-\ln(1 - F)]$ का समाकलन
$C_S = C_0$ से, $F = 0$ पर, आरंभ करने पर $\ln(C_S/C_0) = (1/D - 1)\ln(1 - F)$ प्राप्त होता है, और
$C_L = C_S/D$।
\end{proof}

असंगत तत्त्व के लिए $1/D - 1$ बड़ा और धनात्मक होता है: अवशेष
गलन के आरंभ से $D$ के कुछ गुना के भीतर खाली हो जाता है, और तात्क्षणिक गलित
$C_0/D$ से लगभग शून्य तक गिर जाता है। \omterm{thm:b3:trace-element-modelling:fractional}{भिन्नात्मक गलन} का अवशेष \omterm{thm:b3:trace-element-modelling:batch}{बैच गलन} के अवशेष से कहीं अधिक क्षीण होता है।
गलित एक-एक करके उद्गारित नहीं होते,
फिर भी वे वाहिकाओं में एकत्र होकर जमा होते हैं, और जमा गलित, जिसे नीचे परिभाषित किया गया है,
आकृति में डैश रेखा से दर्शाया गया है।

\begin{omfigure}
\begin{tikzpicture}
\begin{groupplot}[group style={group size=2 by 1, horizontal sep=1.6cm},
  omrecord, width=0.47\linewidth, height=0.44\linewidth, title style={font=\small},
  xlabel={गलन-अंश $F$}, xmin=0, xmax=0.3, ymode=log, ymin=0.01, ymax=150,
  xtick={0,0.1,0.2,0.3}, xticklabels={0,0.1,0.2,0.3},
  ytick={0.01,0.1,1,10,100}, yticklabels={0.01,0.1,1,10,100},
  legend style={font=\scriptsize, legend cell align=left}]
\nextgroupplot[title={(क) बैच (ठोस), समुच्चित (डैश)},
  ylabel={$C_L/C_0$}, legend style={at={(0.97,0.97)}, anchor=north east}]
% computed in figdata/bachelor-3/melting_models.py
\addplot[omDef, thick] table[x=f, y=b001] {figdata/out/bachelor-3-melting_models-curves.dat};
\addlegendentry{$D = 0.01$}
\addplot[red!70!black, thick] table[x=f, y=b01] {figdata/out/bachelor-3-melting_models-curves.dat};
\addlegendentry{$D = 0.1$}
\addplot[black, thick] table[x=f, y=b1] {figdata/out/bachelor-3-melting_models-curves.dat};
\addlegendentry{$D = 1$}
\addplot[omLithLime!70!black, thick] table[x=f, y=b10] {figdata/out/bachelor-3-melting_models-curves.dat};
\addlegendentry{$D = 10$}
\addplot[omDef, thick, densely dashed] table[x=f, y=a001] {figdata/out/bachelor-3-melting_models-curves.dat};
\addplot[red!70!black, thick, densely dashed] table[x=f, y=a01] {figdata/out/bachelor-3-melting_models-curves.dat};
\addplot[omLithLime!70!black, thick, densely dashed] table[x=f, y=a10] {figdata/out/bachelor-3-melting_models-curves.dat};
\nextgroupplot[title={(ख) भिन्नात्मक, तात्क्षणिक}]
\addplot[omDef, thick] table[x=f, y=i001] {figdata/out/bachelor-3-melting_models-curves.dat};
\addplot[red!70!black, thick] table[x=f, y=i01] {figdata/out/bachelor-3-melting_models-curves.dat};
\addplot[black, thick] table[x=f, y=i1] {figdata/out/bachelor-3-melting_models-curves.dat};
\addplot[omLithLime!70!black, thick] table[x=f, y=i10] {figdata/out/bachelor-3-melting_models-curves.dat};
\end{groupplot}
\end{tikzpicture}

\par\smallskip{\small गलित में अनुरेखी तत्त्व का समृद्धीकरण बनाम
गलन-अंश, चार सकल विभाजन गुणांकों के लिए। (क) बैच गलित (ठोस)
और \omterm{def:b3:trace-element-modelling:fractional}{समुच्चित भिन्नात्मक गलित} (डैश; $D = 1$ के लिए दोनों 1 पर बने रहते हैं)। (ख) \omterm{thm:b3:trace-element-modelling:fractional}{भिन्नात्मक गलन} के
तात्क्षणिक गलित: $D = 0.01$ वाला तत्त्व
कुछ प्रतिशत गलन के बाद अवशेष से समाप्त हो जाता है।}\label{fig:b3:trace-element-modelling:curves}
\end{omfigure}

\begin{definition}[समुच्चित भिन्नात्मक गलित]\label{def:b3:trace-element-modelling:fractional}
\emph{समुच्चित भिन्नात्मक गलित}\index{समुच्चित भिन्नात्मक गलित} \omterm{thm:b3:trace-element-modelling:fractional}{भिन्नात्मक गलन} के
$0$ से $F$ तक के सभी तात्क्षणिक गलितों का मिश्रण है,
जो अवशेष के साथ आगे कोई अभिक्रिया किए बिना जमा किया गया हो।
\end{definition}

\begin{proposition}[समुच्चित गलित का संघटन]
\label{prop:b3:trace-element-modelling:aggregate}
स्थिर $D$ के लिए \omterm{def:b3:trace-element-modelling:fractional}{समुच्चित भिन्नात्मक गलित} में निम्न संबंध होता है:
\[
\frac{\bar C_L}{C_0} = \frac{1 - (1 - F)^{1/D}}{F}.
\]
\end{proposition}

\begin{proof}
निकाला गया तत्त्व स्रोत में से अवशेष घटाने पर प्राप्त होता है,
$C_0 - (1 - F)C_S = C_0[1 - (1 - F)^{1/D}]$ \cref{thm:b3:trace-element-modelling:fractional} से,
गलित के द्रव्यमान $F$ पर फैलाया गया है।
\end{proof}

जब $(1 - F)^{1/D}$ छोटा हो, तो समुच्चित गलित में तत्त्व की पूरी मात्रा आ जाती है, और
$\bar C_L \approx C_0/F$ होता है, चाहे $D$ कुछ भी हो: बहुत असंगत तत्त्वों के लिए
\omterm{def:b3:trace-element-modelling:fractional}{समुच्चित भिन्नात्मक गलित} और छोटे $D$ वाले बैच गलित एक-दूसरे से मेल खाते हैं,
और दोनों $F$ को मापते हैं।

वास्तविक गलन इन दोनों अंत्य-सदस्यों के बीच होता है। अवशेष अपने रंध्रों में कुछ गलित रखता है,
और जो गलित निकलता है, वह उस छिद्रता से ऊपर का अधिशेष होता है।

\begin{definition}[गतिक गलन]\label{def:b3:trace-element-modelling:dynamic}
\emph{गतिक गलन}\index{गतिक गलन} में (जिसे सतत गलन भी कहते हैं),
अवशेष अपने रंध्रों में गलित का एक स्थिर द्रव्यमान-अंश $\phi$ रखता है,
जो उसके क्रिस्टलों के साथ साम्य में होता है, और उससे अधिक जितना भी गलित बनता है,
वह निकाल लिया जाता है।
\end{definition}

\begin{proposition}[गतिक गलन]\label{prop:b3:trace-element-modelling:dynamic}
छिद्रता $\phi$ वाला \omterm{def:b3:trace-element-modelling:dynamic}{गतिक गलन}, प्रभावी विभाजन गुणांक $D' = (1 - \phi)D + \phi$ वाला
\omterm{thm:b3:trace-element-modelling:fractional}{भिन्नात्मक गलन} है, जहाँ $F$ निकले गलित का
द्रव्यमान-अंश है:
\[
\frac{C_L}{C_0} = \frac{1}{D'}(1 - F)^{1/D' - 1}.
\]
\end{proposition}

\begin{proof}
अवशेष में क्रिस्टल होते हैं जिनका अंश $1 - \phi$ है, और रंध्र-गलित होता है जिसका अंश $\phi$ है,
दोनों $C_L$ पर गलित के साथ साम्य में हैं, इसलिए उसमें
$C_S = [(1 - \phi)D + \phi]C_L = D'C_L$ होता है।
\cref{thm:b3:trace-element-modelling:fractional} का प्रमाण $D'$ के साथ लागू होता है।
\end{proof}

$\phi = 0$ के लिए मॉडल \omterm{thm:b3:trace-element-modelling:fractional}{भिन्नात्मक गलन} है; $\phi \to 1$ के लिए वह
\omterm{thm:b3:trace-element-modelling:batch}{बैच गलन} की ओर पहुँचता है। एक प्रतिशत छिद्रता पहले गलित के समृद्धीकरण को
$1/\phi = 100$ पर सीमित कर देती है, चाहे $D$ कितना भी छोटा हो:
रंध्र-गलित अवशेष के लिए बफ़र का काम करता है।

\begin{omfigure}
\begin{tikzpicture}[font=\small, >=Stealth]
\tikzset{rock/.style={draw, fill=omLithLime!25, minimum width=2.3cm, minimum height=2.0cm},
  melt/.style={draw, fill=red!25}}
% (क) batch
\node[rock] (a) at (0,0) {};
\foreach \x/\y in {-0.7/0.5, 0.3/0.6, -0.2/-0.1, 0.7/-0.4, -0.6/-0.6, 0.4/0.0}
  \fill[red!55] (\x,\y) circle (0.09);
\node[align=center, font=\scriptsize] at (0,-1.45) {अवशेष $1-F$ पर $DC_L$\\गलित $F$ पर $C_L$};
\draw[->, thick, red!70!black] (0,1.05) -- (0,1.7) node[above, font=\scriptsize] {एक साथ};
\node[font=\small] at (0,-2.25) {(क) बैच};
% (ख) fractional
\node[rock] (b) at (4,0) {};
\draw[->, thick, red!70!black] (4,1.05) -- (4,1.7);
\foreach \y/\s in {1.85/0.06, 2.05/0.08, 2.27/0.1}
  \fill[red!55] (4,\y) circle (\s);
\node[align=center, font=\scriptsize] at (4,-1.45) {हर बूँद बनते ही\\निकल जाती है};
\node[font=\small] at (4,-2.25) {(ख) भिन्नात्मक};
% (ग) dynamic
\node[rock] (c) at (8,0) {};
\foreach \x/\y in {7.4/0.5, 8.4/0.3, 7.8/-0.5, 8.6/-0.6}
  \fill[red!55] (\x,\y) circle (0.06);
\draw[->, thick, red!70!black] (8,1.05) -- (8,1.7);
\foreach \y/\s in {1.85/0.06, 2.05/0.08, 2.27/0.1}
  \fill[red!55] (8,\y) circle (\s);
\node[align=center, font=\scriptsize] at (8,-1.45) {रंध्र $\phi$ रखते हैं;\\अधिशेष निकलता है};
\node[font=\small] at (8,-2.25) {(ग) गतिक};
\end{tikzpicture}

\par\smallskip{\small गलित के गलते स्रोत से बाहर निकलने के तीन तरीके। (क)
बैच: गलित पूरे अवशेष के साथ रहता है, फिर बाहर निकलता है। (ख) भिन्नात्मक:
हर वृद्धि तुरंत बाहर निकल जाती है; संचित वृद्धियाँ समुच्चित
गलित। (ग) गतिक: अवशेष अपने रंध्रों में गलित का एक छोटा द्रव्यमान-अंश $\phi$
रखता है।}\label{fig:b3:trace-element-modelling:modes}
\end{omfigure}

\section{गलन, फिर क्रिस्टलन}\label{sec:b3:trace-element-modelling:crystallisation}

अपने स्रोत से निकला गलित शायद ही कभी अपरिवर्तित रूप में ज्वालामुखी से बाहर आता है। वह
ऊपर जाते समय ओलिवीन, फिर पाइरॉक्सीन और प्लैजियोक्लेज़ क्रिस्टलित करता है, वह दीवार की चट्टान को आत्मसात कर सकता है,
और दूसरे बैचों के साथ मिल भी सकता है। पुस्तक 4 (स्नातक वर्ष 2) के नियम बारी-बारी से लागू होते हैं।
लावा स्रोत के अपने अंश $F$ का गलित है, जो फिर एक ऐसा द्रव है
जिसने भिन्नात्मक क्रिस्टलन के माध्यम से अपने द्रव्यमान का अंश $f$ बनाए रखा है,
जिससे, $D$ वाले तत्त्व के लिए गलन अवशेष में और $D_c$ के लिए
क्रिस्टलित होने वाले खनिजों में, निम्न संबंध प्राप्त होता है:
\[
C_{\text{लावा}} = C_0\,\frac{1}{D + F(1 - D)}\,f^{\,D_c - 1}.
\]
दूसरा गुणक हर उस तत्त्व के लिए समान है जिसे क्रिस्टलित होने वाले
खनिज अस्वीकार करते हैं, जिसके लिए $D_c \approx 0$ है: यह उन सबको $1/f$ से गुणा करता है। ऐसे तत्त्वों के अनुपात
इसलिए स्रोत पर ही तय होते हैं, और जिस लावा से ओलिवीन निकल गया है, वह अपने प्राथमिक गलित का
La/Yb रखता है, जबकि उसका लैंथेनम बढ़ जाता है। सांद्रता से $F$ का प्रतिलोम निकालने
के लिए $f$ भी चाहिए; अनुपात से स्रोत का प्रतिलोम निकालने के लिए नहीं।
दुर्लभ मृदाओं में प्लैजियोक्लेज़ अपवाद है: वह यूरोपियम लेता है,
और उसके निकल जाने से पुस्तक 4 (स्नातक वर्ष 2) की यूरोपियम विसंगति बनी रहती है।

\section{दुर्लभ मृदा प्रतिरूप पढ़ना}\label{sec:b3:trace-element-modelling:patterns}

स्पाइनल लर्ज़ोलाइट का एक गलित और गार्नेट लर्ज़ोलाइट का एक गलित हल्की दुर्लभ मृदाएँ साझा करते हैं,
जिन्हें कोई भी अवशेष नहीं रखता, और भारी मृदाओं में भिन्न होते हैं,
जिन्हें गार्नेट रखता है। कॉन्ड्राइट-प्रसामान्यीकृत आरेख पर गार्नेट गलित
भारी सिरे पर तीव्र ढाल और समतल तल वाले होते हैं; स्पाइनल गलित मंद होते हैं।
$F$ घटने के साथ ढाल बढ़ती है।

\begin{omfigure}
\begin{tikzpicture}
\begin{axis}[omrecord, width=0.8\linewidth, height=0.5\linewidth, ymode=log,
  xmin=56.5, xmax=71.5, ymin=1, ymax=1000, xtick={57,58,60,62,63,64,65,70,71},
  xticklabels={La,Ce,Nd,Sm,Eu,Gd,Tb,Yb,Lu}, ytick={1,10,100,1000}, yticklabels={1,10,100,1000},
  ylabel={नमूना / CI कॉन्ड्राइट},
  legend style={font=\scriptsize, at={(0.98,0.98)}, anchor=north east,
  legend cell align=left, legend columns=2}]
% computed in figdata/bachelor-3/melting_models.py
\addplot[omDef, thick, mark=*, mark size=1.3pt] table[x=z, y=sp01] {figdata/out/bachelor-3-melting_models-ree.dat};
\addlegendentry{स्पाइनल, $F = 1\,\%$}
\addplot[red!70!black, thick, mark=square*, mark size=1.3pt] table[x=z, y=gt01] {figdata/out/bachelor-3-melting_models-ree.dat};
\addlegendentry{गार्नेट, $F = 1\,\%$}
\addplot[omDef, thick, densely dashed, mark=*, mark size=1.3pt] table[x=z, y=sp05] {figdata/out/bachelor-3-melting_models-ree.dat};
\addlegendentry{स्पाइनल, $5\,\%$}
\addplot[red!70!black, thick, densely dashed, mark=square*, mark size=1.3pt] table[x=z, y=gt05] {figdata/out/bachelor-3-melting_models-ree.dat};
\addlegendentry{गार्नेट, $5\,\%$}
\addplot[omDef, thick, densely dotted, mark=*, mark size=1.3pt] table[x=z, y=sp20] {figdata/out/bachelor-3-melting_models-ree.dat};
\addlegendentry{स्पाइनल, $20\,\%$}
\addplot[red!70!black, thick, densely dotted, mark=square*, mark size=1.3pt] table[x=z, y=gt20] {figdata/out/bachelor-3-melting_models-ree.dat};
\addlegendentry{गार्नेट, $20\,\%$}
\addplot[black, thin, mark=none] table[x=z, y=pm] {figdata/out/bachelor-3-melting_models-ree.dat};
\addlegendentry{स्रोत}
\end{axis}
\end{tikzpicture}

\par\smallskip{\small कॉन्ड्राइट-प्रसामान्यीकृत दुर्लभ मृदा प्रतिरूप, बैच गलितों के
आद्य प्रावार स्रोत के, स्पाइनल लर्ज़ोलाइट (नीले वृत्त) और
गार्नेट लर्ज़ोलाइट (लाल वर्ग) से, तीन गलन-अंशों पर; स्रोत
पतली रेखा है। केवल वे नौ दुर्लभ मृदा तत्त्व, जिनके सारणीबद्ध विभाजन गुणांक ज्ञात हैं,
आरेखित किए गए हैं।}\label{fig:b3:trace-element-modelling:ree}
\end{omfigure}
% ledger: budahn91:d-ree, ms95:pm-trace, lodders:ci-table, white13:ex73, rfr25:gkr-5gpa

प्रावार के ऐलुमिनस खनिज के रूप में गार्नेट का स्थान स्पाइनल ले लेता है, लगभग 60
से \qty{90}{km} गहराई पर। जिस गलित की भारी दुर्लभ मृदाएँ समतल और निम्न हों,
और हल्के-से-भारी तीव्र ढाल वाली हो, उसने इसलिए अपने अवशेष में गार्नेट छोड़ा है
और कम से कम उतनी गहराई पर गलना आरंभ किया है।
% ledger: white13:sp-gt

एक प्रतिरूप दो अज्ञात राशियों, $F$ और स्रोत, को मिला देता है।
दो तत्वों को एक-दूसरे के विरुद्ध आलेखित करने से दोनों अलग हो जाते हैं।

\begin{definition}[अनुपात--अनुपात आरेख]\label{def:b3:trace-element-modelling:ratios}
\emph{अनुपात--अनुपात आरेख}\index{अनुपात--अनुपात आरेख} लावाओं के एक समूह के लिए,
अत्यधिक असंगत तत्त्व $H$ का, कम असंगत तत्त्व $M$ से अनुपात,
$H$ की सांद्रता (या किसी अन्य अनुपात) के विरुद्ध आलेखित करता है।
आंशिक गलन, क्रिस्टलन और मिश्रण इस पर भिन्न आकृतियों की रेखाएँ
खींचते हैं।
\end{definition}

\begin{proposition}[बैच गलन रेखा]\label{prop:b3:trace-element-modelling:line}
एक स्रोत के बैच गलितों के लिए, जहाँ $D_H = 0$ और स्थिर $D_M$ है, बिंदु
$(C_H, C_H/C_M)$ सीधी रेखा पर स्थित होते हैं
\[
\frac{C_H}{C_M} = \frac{D_M}{C_0^M}\,C_H + \frac{C_0^H(1 - D_M)}{C_0^M}.
\]
\end{proposition}

\begin{proof}
$D_H = 0$ होने पर $C_H = C_0^H/F$ है, अतः $1/F = C_H/C_0^H$। तब
$C_H/C_M = (C_0^H/F)\,(D_M + F(1 - D_M))/C_0^M = C_0^HD_M/(FC_0^M) +
C_0^H(1 - D_M)/C_0^M$, और पहला पद $(D_M/C_0^M)\,C_H$ है।
\end{proof}

जो खनिज दोनों तत्त्वों को अस्वीकार करते हैं, उनका भिन्नात्मक क्रिस्टलन लावा को
इस आरेख पर क्षैतिज रूप से खिसकाता है, क्योंकि अनुपात नहीं बदलता; गलन इसे
एक ढलवाँ रेखा पर ले जाता है, जिसकी ढाल $D_M$ के साथ बढ़ती है; दो गलितों का मिश्रण
अतिपरवलय खींचता है। गार्नेट वाले स्रोत का $D_{\mathrm{Yb}}$ बड़ा होता है, अतः
La/Yb बनाम La में तीव्र गलन रेखा प्राप्त होती है।

\begin{method}[दो लावाओं का व्युत्क्रमण]\label{met:b3:trace-element-modelling:invert}
(1) जाँचिए कि लावा सह-जनित हैं और $H$, $M$ की तुलना में कहीं अधिक
असंगत है। (2) प्रत्येक लावा को क्रिस्टलन के लिए सुधारिए, या केवल अनुपात का प्रयोग कीजिए।
(3) दो या अधिक लावाओं से होकर रेखा $C_H/C_M = aC_H + b$ का फ़िट कीजिए। (4) तब
$a = D_M/C_0^M$ और $b = C_0^H(1 - D_M)/C_0^M$: $C_0^M$ मान लिया गया हो, तो $D_M =
aC_0^M$ और $C_0^H = bC_0^M/(1 - D_M)$। (5) प्रत्येक लावा का गलन-अंश
$F = C_0^H/C_H$ होता है।
\end{method}

\begin{omfigure}
\begin{tikzpicture}
\begin{axis}[omrecord, width=0.78\linewidth, height=0.52\linewidth,
  xmode=log, ymode=log, xmin=2, xmax=60, ymin=2, ymax=100,
  xlabel={La (\unit{ppm})}, ylabel={La/Yb},
  xtick={2,5,10,20,50}, xticklabels={2,5,10,20,50},
  ytick={2,5,10,20,50,100}, yticklabels={2,5,10,20,50,100},
  legend style={font=\scriptsize, at={(0.03,0.97)}, anchor=north west,
  legend cell align=left}]
% computed in figdata/bachelor-3/melting_models.py
\addplot[omDef, thick] table[x=la_sp, y=lay_sp] {figdata/out/bachelor-3-melting_models-traj.dat};
\addlegendentry{स्पाइनल लर्ज़ोलाइट}
\addplot[red!70!black, thick] table[x=la_gt, y=lay_gt] {figdata/out/bachelor-3-melting_models-traj.dat};
\addlegendentry{गार्नेट लर्ज़ोलाइट}
\addplot[omDef, only marks, mark=|, mark size=3pt, forget plot] table[x=la_sp, y=lay_sp] {figdata/out/bachelor-3-melting_models-ticks.dat};
\addplot[red!70!black, only marks, mark=|, mark size=3pt, forget plot] table[x=la_gt, y=lay_gt] {figdata/out/bachelor-3-melting_models-ticks.dat};
\addplot[black, only marks, mark=diamond*, mark size=2.8pt] table[x=la, y=layb] {figdata/out/bachelor-3-melting_models-lavas.dat};
\addlegendentry{लावा A और B}
\node[font=\scriptsize, anchor=north] at (axis cs:9.6,5.2) {A};
\node[font=\scriptsize, anchor=south east] at (axis cs:15.2,36) {B};
\node[font=\scriptsize, omDef, anchor=north west] at (axis cs:36.5,17.5) {1\,\%};
\node[font=\scriptsize, omDef, anchor=north west] at (axis cs:11.3,6.2) {5\,\%};
\node[font=\scriptsize, omDef, anchor=north west] at (axis cs:3.14,2.55) {20\,\%};
\node[font=\scriptsize, red!70!black, anchor=south east] at (axis cs:29.5,66) {1\,\%};
\node[font=\scriptsize, red!70!black, anchor=south east] at (axis cs:10.5,24) {5\,\%};
\node[font=\scriptsize, red!70!black, anchor=south east] at (axis cs:3.09,7.0) {20\,\%};
\end{axis}
\end{tikzpicture}

\par\smallskip{\small आद्य प्रावार स्रोत के बैच गलित, La/Yb
बनाम La आरेख पर, स्पाइनल और गार्नेट लर्ज़ोलाइट के लिए, चिह्नों के साथ $F = 1$
2, 5, 10 और \qty{20}{\percent} पर, और सप्ताहांत समस्या के दो काल्पनिक लावा,
प्रत्येक लावा केवल एक ही वक्र पर स्थित है।}\label{fig:b3:trace-element-modelling:ratios}
\end{omfigure}

लघुगणकीय अक्षों पर प्रतिज्ञप्ति की सीधी रेखा मुड़ जाती है, और
$D_{\mathrm{La}} \ne 0$ होने पर वह रैखिक अक्षों पर केवल लगभग सीधी होती है; आकृति के वक्र
यथार्थ होते हैं। स्पाइनल स्रोत ऐसा गलित नहीं बना सकता जिसका La/Yb
$(C_0^{\mathrm{La}}/D_{\mathrm{La}})/(C_0^{\mathrm{Yb}}/D_{\mathrm{Yb}})$,
से अधिक हो। आद्य प्रावार के लिए यह मान लगभग 40 है; गार्नेट स्रोत लगभग 120 तक पहुँच सकता है।

\section{अभ्यास}

\begin{exercise}[$\star$]\label{exo:b3:trace-element-modelling:1}
समैरियम के खनिज गुणांक क्लिनोपाइरॉक्सीन में 0.40,
गार्नेट में 0.11 और ओलिवीन, ऑर्थोपाइरॉक्सीन तथा स्पाइनल में 0.016 हैं। इनके आधार पर
स्पाइनल लर्ज़ोलाइट (\qty{12}{\percent} क्लिनोपाइरॉक्सीन) में समैरियम का सकल विभाजन गुणांक
और अध्याय के गार्नेट लर्ज़ोलाइट (\qty{19.5}{\percent} क्लिनोपाइरॉक्सीन,
\qty{16.7}{\percent} गार्नेट) में उसका सकल विभाजन गुणांक परिकलित कीजिए।
\end{exercise}
% ledger: budahn91:d-ree

\begin{exercise}[$\star$]\label{exo:b3:trace-element-modelling:2}
$D = 0.05$ वाला एक तत्त्व किसी स्रोत में \qty{1}{ppm} पर उपस्थित है।
\omterm{thm:b3:trace-element-modelling:batch}{बैच गलन} से बने गलित में उसकी सांद्रता $F = 0.02$ और $0.10$ पर परिकलित कीजिए।
\end{exercise}

\begin{exercise}[$\star$]\label{exo:b3:trace-element-modelling:3}
\omterm{thm:b3:trace-element-modelling:batch}{बैच गलन} से अधिकतम कितना संवर्धन $C_L/C_0$ किसी $D = 0.01$ वाले
तत्त्व में संभव है? और $D = 0$ वाले तत्त्व के लिए $F = \qty{5}{\percent}$ पर कितना?
(नहीं)
\end{exercise}

\begin{exercise}[$\star$]\label{exo:b3:trace-element-modelling:4}
एक लावा में \qty{30}{ppm} La और \qty{2.5}{ppm} Yb है। CI कॉन्ड्राइट के मान
La 0.232 और Yb \qty{0.163}{ppm} लेकर उसका La/Yb और कॉन्ड्राइट-प्रसामान्यीकृत
$(\mathrm{La}/\mathrm{Yb})_N$ परिकलित कीजिए।
\end{exercise}
% ledger: lodders:ci-table

\begin{exercise}[$\star\star$]\label{exo:b3:trace-element-modelling:5}
$D = 0.01$ वाले तत्त्व के लिए \qty{5}{\percent} \omterm{thm:b3:trace-element-modelling:batch}{बैच गलन} और \omterm{thm:b3:trace-element-modelling:fractional}{भिन्नात्मक गलन} के बाद
दोनों के अवशेषों की तुलना कीजिए। उस क्षण तात्क्षणिक भिन्नात्मक गलित में
तत्त्व की सांद्रता कितनी होती है?
\end{exercise}

\begin{exercise}[$\star\star$]\label{exo:b3:trace-element-modelling:6}
अध्याय के हल किए गए उदाहरण में (मोड 58, 27, 12, \qty{3}{\percent};
गलन-अनुपात 20, 25, 45, \qty{10}{\percent} ओलिवीन के लिए,
ऑर्थोपाइरॉक्सीन, क्लिनोपाइरॉक्सीन और स्पाइनल के लिए), लैंथेनम के गुणांक
0.0003, 0.002, 0.053 और 0.01 हैं। $D_0$ और $P$ ज्ञात करके
\omterm{thm:b3:trace-element-modelling:batch}{बैच गलन} से बने गलित में $F = 0.02$ और $0.10$ पर लैंथेनम का संवर्धन, अ-मोडल और मोडल दोनों के लिए, परिकलित कीजिए।
\end{exercise}
% ledger: white13:ex73

\begin{exercise}[$\star\star$]\label{exo:b3:trace-element-modelling:7}
$F = 0.10$ पर, $D = 0.1$ और $D = 0.01$ के लिए, बैच गलित और \omterm{def:b3:trace-element-modelling:fractional}{समुच्चित भिन्नात्मक गलित} की तुलना कीजिए।
कौन-सा अधिक भिन्न है, और क्यों?
\end{exercise}

\begin{exercise}[$\star\star$]\label{exo:b3:trace-element-modelling:8}
$D = 0.001$ वाला एक तत्त्व $\phi = 0.01$ के साथ \omterm{def:b3:trace-element-modelling:dynamic}{गतिक गलन} से गलता है।
उसके $D'$, पहले निकाले गए गलित का संवर्धन और $F = 0.01$ पर निकाले गए गलित का संवर्धन
परिकलित कीजिए, और इनकी तुलना \omterm{thm:b3:trace-element-modelling:fractional}{भिन्नात्मक गलन} से कीजिए।
\end{exercise}

\begin{exercise}[$\star\star$]\label{exo:b3:trace-element-modelling:9}
अध्याय के स्पाइनल लर्ज़ोलाइट में $D_{\mathrm{La}} = 0.0078$ और
$D_{\mathrm{Sm}} = 0.062$ है। $F = 0.01$ और $F = 0.10$ पर \omterm{thm:b3:trace-element-modelling:batch}{बैच गलन} La/Sm को स्रोत के मान से
किस गुणक से बढ़ाता है? इससे
$F$ के अनुमान के लिए अनुपातों के उपयोग के बारे में क्या पता चलता है?
\end{exercise}

\begin{exercise}[$\star\star\star$]\label{exo:b3:trace-element-modelling:10}
दिखाइए कि \omterm{def:b3:trace-element-modelling:fractional}{समुच्चित भिन्नात्मक गलित} 0 से $F$ तक के तात्क्षणिक गलितों का माध्य है,
और यह $C_0/F$ की ओर तब प्रवृत्त होता है जब $(1 - F)^{1/D} \to 0$ हो। $D = 0.01$ के लिए
किस $F$ से यह $C_0/F$ के \qty{1}{\percent} के भीतर है?
\end{exercise}

\begin{exercise}[$\star\star\star$]\label{exo:b3:trace-element-modelling:11}
सह-उत्पन्न लावा P और Q में La 20 और \qty{10}{ppm} है, और La/Yb 10 और 6 है।
मान लीजिए कि लैंथेनम पूर्णतः असंगत तत्त्व है और स्रोत में
\qty{0.44}{ppm} Yb है। अध्याय की विधि लगाकर $D_{\mathrm{Yb}}$,
स्रोत का La और प्रत्येक लावा का गलन-अंश ज्ञात कीजिए।
\end{exercise}

\begin{exercise}[$\star\star\star$]\label{exo:b3:trace-element-modelling:12}
अभ्यास~6 के \omterm{def:b3:trace-element-modelling:nonmodal}{अ-मोडल गलन} में प्रत्येक खनिज किस गलन-अंश पर
समाप्त हो जाता है? लैंथेनम के लिए $F = 0.2$ पर अवशेष का सकल गुणांक परिकलित कीजिए,
और बताइए कि अध्याय का सूत्र पहली समाप्ति के बाद
क्यों लागू नहीं किया जा सकता।
\end{exercise}

\section{समस्या: ऐल्ड्रेन द्वीप के दो लावा}

\begin{problem}[{सप्ताहांत समस्या --- स्रोत, गहराई और दो काल्पनिक बेसाल्टों के गलन की मात्रा,
उनके लैंथेनम और इटरबियम के मानों से}]\label{pb:b3:trace-element-modelling:1}
ऐल्ड्रेन द्वीप पर, जो एक काल्पनिक ज्वालामुखीय द्वीप है, लावा A में \qty{9.6}{ppm} La
और \qty{1.7}{ppm} Yb है, और लावा B में \qty{15.6}{ppm} La और \qty{0.45}{ppm} Yb है।
दोनों को एक आद्य प्रावार स्रोत के प्राथमिक गलित माना जाता है, जिसमें
\qty{0.648}{ppm} La और \qty{0.441}{ppm} Yb है (CI कॉन्ड्राइट: La 0.232, Yb
\qty{0.163}{ppm})। खनिज गुणांक: La क्लिनोपाइरॉक्सीन में 0.058 और 0.001
गार्नेट में तथा ओलिवीन, ऑर्थोपाइरॉक्सीन और स्पाइनल में; Yb क्लिनोपाइरॉक्सीन में 0.58,
गार्नेट में 4.5 और ओलिवीन, ऑर्थोपाइरॉक्सीन और स्पाइनल में 0.16 हैं।
स्पाइनल लर्ज़ोलाइट में \qty{58}{\percent} ओलिवीन, \qty{27}{\percent} ऑर्थोपाइरॉक्सीन, \qty{12}{\percent}
क्लिनोपाइरॉक्सीन, \qty{3}{\percent} स्पाइनल; गार्नेट लर्ज़ोलाइट में \qty{54.1}{\percent} ओलिवीन, \qty{9.8}{\percent}
ऑर्थोपाइरॉक्सीन, \qty{19.5}{\percent} क्लिनोपाइरॉक्सीन, \qty{16.7}{\percent} गार्नेट। गलन मोडल है।
% ledger: ms95:pm-trace, lodders:ci-table, budahn91:d-ree, white13:ex73, rfr25:gkr-5gpa

\medskip\noindent\textbf{भाग I --- दो स्रोत।}
\begin{enumerate}
\item स्पाइनल लर्ज़ोलाइट के $D_{\mathrm{La}}$ और $D_{\mathrm{Yb}}$ परिकलित कीजिए।
\item गार्नेट लर्ज़ोलाइट के लिए इन्हें परिकलित कीजिए।
शेष समस्या के लिए चारों मानों को दो सार्थक अंकों तक पूर्णांकित कीजिए।
\item दोनों स्रोतों में कौन-सा गुणांक सबसे अधिक भिन्न है, और क्यों?
क्या गार्नेट लर्ज़ोलाइट में इटरबियम संगत तत्त्व है?
\item स्रोत का कॉन्ड्राइट-प्रसामान्यीकृत $(\mathrm{La}/\mathrm{Yb})_N$ परिकलित कीजिए।
इसका प्रतिरूप कैसा दिखता है?
\end{enumerate}

\medskip\noindent\textbf{भाग II --- लावा A का बैच गलन।}
\begin{enumerate}[resume]
\item स्पाइनल लर्ज़ोलाइट के बैच गलितों में $F = 1$, 5 और
\qty{10}{\percent} पर La परिकलित कीजिए।
\item वह गलन-अंश ज्ञात कीजिए जो लावा A का लैंथेनम देता है।
\item इसका इटरबियम और La/Yb पूर्वानुमानित कीजिए, और इनकी तुलना लावा A से कीजिए।
\item अवशेष में La और Yb, तथा स्रोत के
लैंथेनम का वह अंश जो उसमें शेष रहता है, परिकलित कीजिए।
\item यदि लैंथेनम को पूर्णतः असंगत मान लिया जाए, तो कौन-सा गलन-अंश
प्राप्त होगा?
\end{enumerate}

\medskip\noindent\textbf{भाग III --- लावा A का भिन्नात्मक गलन।}
\begin{enumerate}[resume]
\item तात्कालिक भिन्नात्मक गलित में $F = 0.1$, 1 और
\qty{6}{\percent} पर La परिकलित कीजिए।
\item $F = \qty{6.0}{\percent}$ पर \omterm{def:b3:trace-element-modelling:fractional}{समुच्चित भिन्नात्मक गलित} में La परिकलित कीजिए।
\item वह गलन-अंश ज्ञात कीजिए जिस पर समुच्चित गलित में लावा A का
लैंथेनम प्राप्त होता है। यह प्रश्न~9 के उत्तर के इतना निकट क्यों है?
\item अवशेष के लैंथेनम की तुलना प्रश्न~8 के लैंथेनम से कीजिए।
\item $F$ का कौन-सा अनुमान आप रखते हैं?
\end{enumerate}

\medskip\noindent\textbf{भाग IV --- लावा B और गार्नेट।}
\begin{enumerate}[resume]
\item दोनों लावाओं का La/Yb और $(\mathrm{La}/\mathrm{Yb})_N$ परिकलित कीजिए।
\item स्पाइनल लर्ज़ोलाइट का बैच गलित अधिकतम कितना La/Yb
प्राप्त कर सकता है?
\item किस गलन-अंश पर स्पाइनल लर्ज़ोलाइट के गलित का La/Yb
लावा B के बराबर होगा, और उसमें कितना लैंथेनम होगा? निष्कर्ष निकालिए।
\item गार्नेट लर्ज़ोलाइट के लिए, लावा B के लैंथेनम से उसका गलन-अंश
ज्ञात कीजिए, और उसका इटरबियम पूर्वानुमानित कीजिए।
\item दोनों स्रोत उन गहराइयों के बारे में क्या कहते हैं जिन पर A और B
का गलन आरंभ हुआ? क्या लावा A गार्नेट लर्ज़ोलाइट से आ सकता है?
\item गार्नेट स्रोत की \omterm{thm:b3:trace-element-modelling:batch}{बैच गलन} रेखा La/Yb $= a\,
\mathrm{La} + b$ का ढलान और अंतःखंड, $D_{\mathrm{La}}$ को शून्य मानकर, परिकलित कीजिए
और उसके विरुद्ध लावा B की जाँच कीजिए।
\item यदि उद्गार से पहले लावा A ने ओलिवीन के रूप में अपने द्रव्यमान का
\qty{10}{\percent} खो दिया होता, तो उसका La/Yb और गलन-अंश कैसे बदलते?
\item इन दोनों लावाओं के लिए अध्याय के आरंभ में पूछे गए प्रश्न का उत्तर दीजिए।
\item लावा A के गलन की मात्रा बताइए।
\end{enumerate}
\end{problem}
